\documentclass[11pt]{article}

\usepackage[a4paper, margin=2.2cm]{geometry}
\usepackage[T1]{fontenc}
\usepackage{lmodern}
\usepackage{booktabs}
\usepackage{enumitem}
\usepackage{hyperref}
\usepackage{microtype}

\setlist{itemsep=2pt, topsep=4pt}
\setlength{\parskip}{4pt}
\setlength{\parindent}{0pt}

\title{Few-Shot Classification with DINOv2 Embeddings\\
\large Project Plan}
\author{
    \texttt{m.mestnikov@innopolis.university}\\
    \texttt{e.moskal@innopolis.university}\\
    \texttt{k.sadykov@innopolis.university}\\
    \texttt{m.klimov@innopolis.university}
}
\date{October 2026}

\begin{document}

\maketitle

\section{Project Scope}

We classify the 37 cat and dog breeds of Oxford-IIIT Pet when only
$K \in \{1, 5, 10, 20\}$ labeled images per breed are available for training.

\textbf{Research question.} With the same $K$ labeled images per breed, how much
better do frozen self-supervised DINOv2 embeddings classify breeds than a
supervised ImageNet ResNet-50 transferred to the task, and which breeds remain
hard for both?

The main condition is a frozen DINOv2 ViT-S/14 encoder with a lightweight
classifier trained on its embeddings. We use ViT-S/14 because, with 21M parameters,
it is close in size to ResNet-50 (25M), so the comparison is not decided by model
capacity; it also runs fast enough on a laptop. The baseline is a ResNet-50 pretrained with
labels on ImageNet and transferred to our task in three ways. Each way replaces the 1000-class ImageNet head with a new
37-class head and differs only in what else is trained: nothing, the last residual
stage, or the whole network.
Both conditions see the same images, the same splits and the same support sets.

\section{Dataset Preparation}

\textbf{Dataset.} Oxford-IIIT Pet: 7,349 images of 37 breeds (12 cat, 25 dog),
about 200 images per breed. We load it with \texttt{torchvision.datasets.OxfordIIITPet}
and use only the breed labels.

\textbf{Split.}
\begin{itemize}
    \item \emph{Test}: the official test split (3,669 images). Used once, for final metrics.
    \item \emph{Validation}: 20 images per breed (740 images) taken from the
          official trainval split with a fixed seed. Used to choose
          hyperparameters and the number of fine-tuning epochs.
    \item \emph{Train pool}: the remaining trainval images (about 2,940,
          at least 70 per breed). Few-shot support sets are drawn only from here.
\end{itemize}

\textbf{Support sets.} For each $K$ we draw 5 support sets with seeds 0--4.
Every method trains on exactly the same support sets, so results are paired
across methods. Each test image belongs to the test split only, which rules out
train/test leakage on our side.

\textbf{Preprocessing.} Resize the shorter side to 256, center crop to
$224 \times 224$, normalize with ImageNet mean and standard deviation.
224 is a multiple of the DINOv2 patch size (14), so both backbones get the same input.
Fine-tuning runs add random resized crops and horizontal flips to the support images.

\section{Experimental Plan}

\subsection{Models}

\smallskip
\begin{tabular}{@{}lll@{}}
\toprule
Backbone & Pretraining & Embedding \\
\midrule
DINOv2 ViT-S/14 & self-supervised, LVD-142M & CLS token, 384-d \\
DINOv2 ViT-B/14 (ablation) & self-supervised, LVD-142M & CLS token, 768-d \\
ResNet-50 (\texttt{IMAGENET1K\_V2}) & supervised, ImageNet-1k & global avg.\ pool, 2048-d \\
\bottomrule
\end{tabular}
\smallskip

\subsection{Classifiers on frozen embeddings}

Embeddings are L2-normalized.
\begin{itemize}
    \item \textbf{Prototype (nearest class mean)}: each breed is the mean of its
          $K$ support embeddings; a test image gets the breed with the highest cosine similarity.
    \item \textbf{k-NN}: cosine distance, $k$ chosen on validation ($k=1$ when $K=1$).
    \item \textbf{Logistic regression}: multinomial, L2 penalty, $C$ chosen on validation.
\end{itemize}

\subsection{Supervised transfer-learning baseline}

\begin{enumerate}
    \item \textbf{New head only}: frozen ResNet-50, same three classifiers as above.
          This is the main baseline, because only the representation differs from the DINOv2 condition.
    \item \textbf{New head + \texttt{layer4}}: the new linear head and the last residual stage
          are trained with AdamW and cross-entropy; earlier layers stay frozen.
    \item \textbf{New head + full backbone}: all weights are trained, with a lower learning
          rate for the backbone than for the head.
\end{enumerate}
For variants 2 and 3, BatchNorm layers stay in eval mode (support batches are
too small for reliable batch statistics), and the epoch count is chosen on validation.
We expect variant 3 to overfit at small $K$ and will report it either way.

\subsection{Experiment order}

Each experiment changes one factor relative to an earlier one. Every experiment
is run for all four values of $K$ and all five support sets, so each one yields an
accuracy-vs-$K$ curve that we compare with the curve of the earlier experiment.

\smallskip
\begin{tabular}{@{}lll@{}}
\toprule
ID & Setting & Changed vs. \\
\midrule
E1 & ResNet-50 frozen + logistic regression & baseline \\
E2 & DINOv2 ViT-S/14 frozen + logistic regression & E1: backbone \\
E3 & DINOv2 ViT-S/14 + prototype / k-NN & E2: classifier \\
E4 & ResNet-50 + prototype / k-NN & E1: classifier \\
E5 & ResNet-50 new head + \texttt{layer4} fine-tuning & E1: trainable layers \\
E6 & ResNet-50 full fine-tuning & E5: trainable layers \\
E7 & DINOv2 ViT-B/14 + best classifier from E3 & E2/E3: model size \\
\bottomrule
\end{tabular}
\smallskip

E3, E4 and E7 are our ablations: they keep everything fixed except the classifier
(E3, E4) or the DINOv2 model size (E7), which shows how much of the result comes from
the classifier choice and how much from encoder capacity.
Every run writes one row (experiment, backbone, classifier, $K$, seed, accuracy,
macro-F1, training time, device) to a CSV file in \texttt{results/}, so every
number in the final report can be regenerated from the code.

\subsection{Metrics}

\textbf{Quality metrics.}
\begin{itemize}
    \item \textbf{Top-1 accuracy} and \textbf{macro-F1} on the test split,
          reported as mean $\pm$ std over the 5 support sets for each $K$.
    \item \textbf{Per-breed accuracy} and the \textbf{confusion matrix} at $K=20$,
          to find the hardest breeds and the most frequent confusions.
\end{itemize}

\textbf{Efficiency metrics.}
\begin{itemize}
    \item \textbf{Training time} per run: fitting a classifier on cached
          embeddings, or fine-tuning the ResNet-50 variants.
    \item \textbf{Parameters} (millions): total and trainable, for each backbone and variant.
    \item \textbf{GFLOPs} per $224 \times 224$ image for one forward pass of each backbone.
\end{itemize}

\subsection{Failure analysis}

We take the breeds with the lowest per-breed accuracy and check whether they are
hard for both backbones or for only one. For at least three misclassified test
images we show the image, the predicted breed and its nearest support images,
and check concrete causes: visually similar breeds, unusual pose or occlusion,
several animals in one image, cluttered background, low resolution.

\section{Implementation Plan}

\textbf{Pipeline.}
\begin{enumerate}
    \item Download the dataset, build the validation split and the train pool,
          sample and save the support-set indices for every $K$ and seed.
    \item Run each frozen backbone once over all 7,349 images and cache the
          embeddings as \texttt{.npy} files.
    \item Train and evaluate the three lightweight classifiers on cached embeddings (CPU, seconds per run).
    \item Fine-tune ResNet-50 variants 2 and 3 on GPU for every $K$ and seed (40 runs).
    \item Aggregate metrics into CSV tables, plot accuracy vs.\ $K$ for all methods,
          compute per-breed accuracy and confusion matrices, collect failure cases.
\end{enumerate}

\textbf{Code layout.} Reusable code in \texttt{src/} (\texttt{data.py},
\texttt{features.py}, \texttt{classifiers.py}, \texttt{finetune.py},
\texttt{evaluate.py}); experiments, plots and the demo in \texttt{notebooks/};
metrics and figures in \texttt{results/}.

\textbf{Libraries and weights.} PyTorch, torchvision, scikit-learn, NumPy, pandas,
Matplotlib. DINOv2 weights from \texttt{torch.hub} (\texttt{facebookresearch/dinov2}),
ResNet-50 weights from torchvision.

\textbf{Hardware.} Apple Silicon (MPS), RTX 4060 Laptop GPU and Kaggle GPUs for fine-tuning.

\textbf{Known limitation.} DINOv2's LVD-142M pretraining set was assembled by
retrieving web images similar to public datasets, so it may contain pet images
close to ours. We cannot check this and will state it in the final report.

\section{Team Responsibilities}

\smallskip
\begin{tabular}{@{}lp{0.72\textwidth}@{}}
\toprule
Member & Responsibility \\
\midrule
Marat & Supervised baseline: ResNet-50 embeddings, fine-tuning variants 2 and 3 on GPU. \\
Evgeni & DINOv2: embedding extraction and caching, prototype / k-NN /
         logistic regression classifiers, ViT-B/14 ablation. \\
Khasan & Data pipeline: download, splits, support-set sampling, preprocessing;
        README, \texttt{requirements.txt}, reproducibility check. \\
Martin & Evaluation: metrics, result tables and plots, per-breed analysis,
         failure cases. \\
\bottomrule
\end{tabular}
\smallskip

All four members write the two-page summary and prepare the demo together, and each
member runs the full pipeline once before submission.

\end{document}
